\chapterimage{container.jpg}

\chapter{¿Esto es tuyo?}\label{sec:packets}

La información viaja por este grafo de redes en forma de paquetes, que contienen una pequeña cantidad de 
información y de meta-información. 

El formato de los paquetes debe estar acordado por ambos extremos de la comunicación, de modo que esos paquetes puedan 
producirse e interpretarse de forma eficiente y automática.

Cuando los paquetes viajan por una o varias redes, pueden perderse o desordenarse. Si tiene que fluir un flujo continuo de datos 
entre dos ordenadores, los paquetes deben contener mecanismos para detectar pérdidas y restaurar el orden correcto.

\section{Tengo un paquete para ti}\label{sec:packets:lan_wan}

No es viable establecer conexiones físicas entre cada par de ordenadores que quieren comunicarse.
En su lugar, se comparten relativamente pocas conexiones para transportar mensajes entre muchos pares de dispositivos.

Si un único dispositivo transmite datos de forma continua durante mucho tiempo, esa conexión queda bloqueada y puede convertirse en un cuello de botella que 
impida comunicarse a otros dispositivos. Para evitar este problema, las conexiones limitan la cantidad máxima de datos
que pueden enviarse sin interrupción. Como resultado, los dispositivos se ven obligados a enviar los datos en ráfagas discretas llamadas 
\conceptRef{packet}{paquetes}. El intercambio de este tipo de mensajes a través de una o varias redes se denomina
\conceptRef{packet switching}{conmutación de paquetes}.

\begin{remark}
% The concept of \concept{packet switching} (crucial to the creation of the Internet)
% was developed in the 1960's (during the cold war) by the US military
% to have a communications network that would work in case of a nuclear strike.
%  
% The problem with the one they had was that, if one station was damaged, the whole system stopped working.
% However, by using packet switching, any connected subgraph of the network would keep working.
% 
% The reasoning for creating a packet-switched network was that they could still retaliate their enemy in case of an attack, which would
% make them a less attractive target.

Leonard Kleinrock, uno de los pioneros de la interconexión de redes, disputa la autoría de la conmutación de paquetes,
que a menudo se atribuye a Paul Baran y a Donald Davies.
\end{remark}

\conceptRef{packet}{Los paquetes} deben ser pequeños para que las conexiones no queden bloqueadas demasiado tiempo. 
Al mismo tiempo, queremos que los paquetes contengan tantos datos como sea posible, porque cada paquete 
contiene datos de \concept{overhead} (sobrecarga), como las direcciones que se comentan en la Sección~\ref{sec:where}, que hay
que enviar además de los datos de payload (carga útil) del usuario. Las redes fijan la longitud máxima de paquete
mediante un parámetro llamado \conceptRef{MTU}{unidad máxima de transmisión} (MTU), \eg\ 1500~bytes en Ethernet.


La transmisión de paquetes a través de un \conceptRef{data link}{enlace de datos} debe hacerse con cuidado,
en particular cuando ese enlace es un \concept{bus} compartido por varios dispositivos.
Cada paquete debe poder distinguirse individualmente del resto, 
lo que puede conseguirse mediante al menos tres estrategias:
\begin{enumerate}
\item \textit{Longitud fija}: la longitud de todos los paquetes es la misma. El emisor y el receptor deben haber
acordado este valor de antemano.

\item \textit{Longitud explícita}: la longitud del paquete se incluye en el propio paquete, \eg\ 
usando el primer byte del paquete. El emisor y el receptor deben ponerse de acuerdo sobre cómo se incluye
esta información y cómo interpretarla.

\item \textit{Secuencias de escape}: ciertas \conceptRef{escape sequence}{secuencias de escape} binarias 
dentro de los datos tienen un significado especial. 
Por ejemplo, podría definirse que el byte \otherBase{11110000} (\otherBase{0xf0}) 
significa ``fin de paquete'': estos bits deben aparecer después de cada paquete, y solo entonces.
% 
Así, el emisor podría empezar a transmitir el contenido de un paquete, seguido de \otherBase{0xf0}.
Ambas partes deben acordar qué secuencias de escape usar, qué significan y cómo 
enviar datos iguales a esas secuencias (\eg\ debería ser posible enviar \otherBase{0xf0f0f0}
sin provocar un ``fin de paquete'' hasta que se hayan transmitido todos los bytes).
\end{enumerate}

\begin{remark}
También es posible señalar el \textit{principio} de un paquete. En ese caso se utiliza una secuencia previa llamada \conceptRef{preamble}{preámbulo}.
\end{remark}


\begin{exercise}
Todas las estrategias descritas arriba se usan hoy en día en escenarios reales, dependiendo de las 
características, los costes y los compromisos de cada una.

Discute cuál de estas estrategias sería la más apropiada para cada uno de estos escenarios:
\begin{itemize}
\item Un único fabricante diseña el hardware y el software que comunica dos extremos.
Una única \conceptRef{data line}{línea de datos} debe \conceptRef{multiplex}{multiplexarse} para varios comandos de control
y también para varios flujos de datos. La prioridad es un uso eficiente de la energía y del buffer (memoria intermedia).

\item Varios dispositivos comparten un \concept{bus}. De vez en cuando envían mensajes de distintas longitudes,
pero no un volumen enorme de datos. La prioridad es el coste.

\item Varios dispositivos comparten un \concept{bus}. Envían continuamente mensajes de distintas longitudes,
intentando maximizar la tasa efectiva de transmisión por el \conceptRef{channel}{canal}. 
La prioridad es el \concept{throughput} (caudal efectivo).
\end{itemize}
\end{exercise}

\section{Por favor, rellena el formulario}\label{sec:packets:format}

Independientemente de la estrategia empleada, los paquetes de datos suelen constar de dos partes: 
\begin{itemize}
 \item \concept{payload} de datos: la información que hay que transmitir, y
 \item \conceptRef{metadata}{metadatos}: la meta-información necesaria para entregar el payload.
\end{itemize}
La longitud puede ser uno de los \conceptRef{field (packet)}{campos} de metadatos incluidos en el paquete, junto con 
otros que ayudan a identificar su tipo y su función.
% 
La ubicación del payload y de los metadatos, así como los campos de metadatos exactos que se incluyen, su longitud y su significado
deben ser conocidos por ambos extremos de la comunicación. Todas estas decisiones constituyen el \conceptRef{packet format}{formato de paquete}
o \conceptRef{packet structure}{estructura de paquete}.

Una de las formas más habituales de disponer el payload y los metadatos es con un formato \conceptRef{header}{cabecera}-\conceptRef{body}{cuerpo}.
En este caso, toda la meta-información se incluye primero, seguida de los datos. 
Otros formatos pueden incluir un \concept{trailer} (cola), también conocido como \concept{tail}, que se usa sobre todo para la \conceptRef{error detection}{detección de errores}
y su corrección, en combinación con la cabecera o en su lugar.
Todas las configuraciones siguientes son posibles.

\begin{minipage}{\linewidth}
\begin{center}
\begin{bytefield}{32}
\bitbox{11}{Header} & \bitbox{21}{Body}\\[-0.35cm]
\end{bytefield}
% 
\begin{bytefield}{32}
\bitbox{6}{Header} & \bitbox{19}{Body} & \bitbox{7}{Trailer/Tail}\\[-0.35cm]
\end{bytefield}
% 
\begin{bytefield}{32}
\bitbox{21}{Body} & \bitbox{11}{Trailer/Tail}\\[-0.35cm]
\end{bytefield}
% 
\begin{bytefield}{32}
\bitbox{32}{Header}
\end{bytefield}
\end{center}
\end{minipage}


El formato del payload de un paquete normalmente lo define el usuario. El formato de la cabecera, en cambio,
está estrictamente definido para que pueda producirse y \conceptRef{parsing}{analizarse} con facilidad.
Este formato suele presentarse mediante un diagrama que ayuda a identificar las posiciones individuales de bit y de byte,
como en el siguiente ejemplo (ficticio):\\

\begin{center}
\begin{bytefield}{32}
\bitheader{0-31} \\
\begin{rightwordgroup}{Header}
\bitbox{16}{Length} & \bitbox{8}{Importance} & \bitbox{4}{Type} & \bitbox{4}{Mode}\\
\bitbox{24}{Destination Address} & \bitbox{8}{(Reserved)} 
\end{rightwordgroup} \\
\begin{rightwordgroup}{Body}
\wordbox[lrt]{1}{Payload} \\
\skippedwords \\
\wordbox[lrb]{1}{} 
\end{rightwordgroup}
\end{bytefield}
\end{center}

\begin{remark}
Los campos también pueden tener longitudes fijas o variables, que no encajan necesariamente con las 
\conceptRef{byte boundary}{fronteras de byte}. Los campos también pueden tener una longitud igual a $1$~bit,
en cuyo caso normalmente se le llama \concept{flag} (indicador) o bit de flag.
\end{remark}


\begin{remark}
En diagramas como el anterior se usan a menudo varias líneas (filas) para que puedan
imprimirse y leerse con facilidad. En ese caso, el significado es el mismo que si todos los campos se hubieran
presentado a lo largo de la misma línea (fila).
\end{remark}

\begin{exercise}
El siguiente contenido de un paquete fue \conceptRef{sniff}{capturado} de una \conceptRef{network}{red},
expresado en \concept{hexadecimal}. Suponiendo que el paquete tiene el formato del ejemplo anterior:
% 
\begin{itemize}
\item Indica el valor de cada campo (longitud, importancia, tipo, modo, dirección de destino y payload.
\item ¿Puedes adivinar el significado del payload?
\end{itemize}
% 
\begin{center}
\otherBase{00 0c 00 f3 bb bb bb 00 68 65 6c 70}
\end{center}
\end{exercise}

\begin{exercise}
El siguiente código acepta dos argumentos: (1, \inlineCode{sys.argv[1]}) la \conceptRef{path}{ruta} de un fichero de entrada, 
y (2, \inlineCode{sys.argv[2]}) la \conceptRef{path}{ruta} de un fichero de salida.
Se leen los primeros 255 bytes del contenido del fichero de entrada, 
se formatean como un paquete y los bytes de ese paquete
se escriben en el fichero de salida.

\begin{center}
\showCode{snippets/simplepacketoutput.py}
\end{center}

\begin{itemize}
\item Describe el \conceptRef{packet format}{formato de paquete} usado en el código, y sus limitaciones.

\item Extiende el formato de paquete para que 
 \begin{itemize}
   \item los paquetes puedan ser más largos de 255 bytes
   \item pueda codificarse la \conceptRef{address}{dirección} del dispositivo de destino
  \end{itemize}
  
\item Proporciona un diagrama de bytes del nuevo formato, 
así como una explicación del sistema de direccionamiento que estés usando.

\item Modifica el código para que produzca paquetes con el formato que has propuesto.
\end{itemize}
\end{exercise}

\begin{remark}
Lenguajes como Python pueden distinguir entre ficheros que contienen texto y ficheros que contienen datos binarios.
En el código anterior, los ficheros se abren en \conceptRef{binary mode}{modo binario} usando los modos \inlineCode{'rb'} y \inlineCode{'wb'}.

Cuando se abre en modo binario, los bytes del fichero se representan directamente mediante un objeto \inlineCode{bytes}, un array de 
enteros entre $0$ y $255$. En \conceptRef{text mode}{modo texto}, esos bytes son procesados (\conceptRef{decode}{decodificados}) además por el método \inlineCode{open},
que devuelve una cadena que puede contener caracteres especiales como acentos o glifos no latinos. \readMore{https://docs.python.org/3.12/library/functions.html\#open}
\end{remark}

\section{No paran de llegar}\label{sec:packets:stream}

Al desarrollar aplicaciones que usan capacidades de red, a menudo resulta útil 
imaginar la comunicación como un \conceptRef{stream}{flujo}, \ie, como una tubería conceptual donde ponemos 
nuestros datos (cualquier cantidad de datos) por un extremo y salen por el otro. Si tenemos
esta capacidad, resulta mucho más fácil enviar ficheros de cualquier tamaño y transmitir una cantidad interminable
de audio o de vídeo. Por desgracia, lo único que tenemos para simular esos flujos son paquetes.

\conceptRef{packet}{Los paquetes} a menudo tienen que reenviarse a través de varias redes. Los \conceptRef{router}{routers} (encaminadores) no son perfectos
y a veces están fuera de servicio o mal configurados, así que no todos los paquetes llegan a su destino. Además,
paquetes distintos pueden entregarse por rutas distintas que pueden tener longitudes distintas, de modo que los paquetes pueden llegar
desordenados.

Si queremos simular un \conceptRef{stream}{flujo} de datos, los paquetes deben contener suficiente meta-información para que puedan detectarse las pérdidas y 
restaurarse el orden correcto. Esto introduce un \concept{overhead} que no es adecuado en todos los escenarios --\eg\ latencia muy baja--,
así que algunas aplicaciones basan su comunicación en red en \conceptRef{message}{mensajes} en lugar de en flujos.

Cuando se requieren flujos, se usa el concepto de \concept{offset} (desplazamiento) para reensamblar varios mensajes en 
un único flujo. Cada paquete contiene algunos bytes del flujo de datos, y el offset indica la posición exacta dentro de ese flujo.

\begin{exercise}

El siguiente diagrama describe un flujo de 48 bytes producido por un nodo origen, 
así como los paquetes en los que se dividió antes de la transmisión:\\[-0.25cm]

\begin{center}
\begin{bytefield}[bitheight=2cm]{48}
\bitheader{0-47} \\
\bitbox{10}{\underline{Packet 0} \\Offset 0 \\Length 10}
\bitbox{14}{\underline{Packet 1} \\Offset 10\\Length 14}
\bitbox{6 }{\underline{Packet 2} \\Offset 24\\Length 6 }
\bitbox{18}{\underline{Packet 3} \\Offset 30\\Length 18}
\end{bytefield}
\end{center}

\begin{itemize}
\item Describe la cabecera mínima que puede usarse para transmitir este flujo.
\item Suponiendo que el receptor recibe primero el Paquete 2, después el Paquete 0, y que los demás paquetes se pierden: ¿qué bytes del flujo conoce el 
  nodo de destino y qué bytes desconoce?
\item ¿Cómo puede saber el nodo origen que algunos paquetes no se entregaron?
\item ¿Qué ocurriría si el destino recibiera un mensaje espurio, con offset de $7$~bytes y longitud de $4$~bytes?
\end{itemize}

\end{exercise}



